%% PRELIMINARIES

%% (2) background and preliminaries
\subsection{Preliminaries}
%
Following the notation in \citep{vaswaniAttentionAllYou2017, radfordImprovingLanguageUnderstanding2018}, we define a Transformer architecture as,
%
\input{ICML2025/math/statements/def-transformer-model}
%
Furthermore, we use the following definition of a bilinear form,
\input{ICML2025/math/statements/def-bilinear-form}
%
Finally, we provide the following definition of autoregressive and bidirectional training objectives,
%
\begin{definition}
\label{def-objective-functions}
%
(\textbf{\emph{Autoregressive and bidirectional objectives}}).
%
Let $U = \{t_1, \dots, t_N\}$ a sequence of tokens.
%
The joint probability of $U$ is factorized autoregressively as follows,
%
\begin{equation}
    \Pr[U] = \Pr[t_1,\dots,t_N] = \Pi_{i=1}^N \Pr[t_i|t_1,\dots,t_{i-1}]\,.
\end{equation}
%
During autoregressive training, a model with a set of parameters $\mathcal{W}$ is optimized to minimize the following negative log-likelihood
%
\begin{equation}
\mathcal{L}(U; \mathcal{W}) = - \sum_{i=1}^{N} \log p_{\mathcal{W}}(t_i \,|\, \{t_j\} : j < i)\,,
\end{equation}
%
where the conditional probabilities are modeled with learnable parameters $\mathcal{W}$.
%
% The join probability of $U$ can also be factorized bidirectionally as follows,
% %
% \begin{equation}
%     \Pr[U] = \Pr[t_1,\dots,t_N] = \Pi_{i=1}^N \Pr[t_i|t_1,\dots,t_{i-1}, t_{i+1}, \dots, t_N]\,.
% \end{equation}
%
During bidirectional training, a model with a set of parameters $\mathcal{W}$ is optimized to minimize the following negative log-likelihood
%
\begin{equation}
\mathcal{L}(U; \mathcal{W}) = - \sum_{i=1}^{N} \log p_{\mathcal{W}}(t_i \,|\, \{t_j\} : j \neq i)\,.
\end{equation}
%
In practice, only a subset $M\in [1,\dots,N]$ of the tokens are predicted as in  Masked Language Modelling (MLM) \citep[see ][]{devlinBERTPretrainingDeep2019,warnerSmarterBetterFaster2024}, leading to the following negative log-likelihood
%
\begin{equation}
\mathcal{L}(U; \mathcal{W}) = - \sum_{i\in M} \log p_{\mathcal{W}}(t_i \,|\, \{t_j\} : j \notin M)\,.
\end{equation}
%
\end{definition}





%% RELATED REMARKS TO SELF-ATTENTION AND BILINEAR FORM
\subsection{Related remarks to Section \ref{sec:self-attention-bilinear-form}}
\label{supp-math-self-attention-bilinear-form}
%
\input{ICML2025-rebuttal/math/statements/remark-self-attention-bilinear-form-projection}
\input{ICML2025-rebuttal/math/statements/remark-self-attention-bilinear-form-multihead}




%% REMARKS ON OBJECTIVE FUNCTION, AUTOREGRESSIVE and BIDIRECTIONAL TRAINING, and BERT MLM
\subsection{Related remarks to Section \ref{sec:self-attention-derivation-gradient}}
\label{supp-math-self-attention-derivation-gradients}
%
Let \( \mathcal{V} = \{t_1, \dots, t_V\} \) be a vocabulary of tokens, and let \( D \) be a dataset of sequences \( U_u \in \mathcal{V}^N \) of length \( N \), where $D = \left\{ U_u = (x^u_1, \dots, x^u_N) \right\}_{u=1}^D \subset \mathcal{V}^N$.
%
For a single sequence \( U_u \) with joint distribution \( p(U_u) \), the chain rule implies that for any permutation \( \sigma \) of \( \{1, \dots, N\} \), the joint probability can be written as,
\begin{equation}
p(U_u) = \prod_{t=1}^{N} p\big(x^u_{\sigma(t)} \mid x^u_{\sigma(1)}, \dots, x^u_{\sigma(t-1)}\big).
\label{eq:chain}
\end{equation}
%
Given a set of conditional parents \( \mathcal{D}_i \subset \{1, \dots, N\} \setminus \{i\} \), define the **local negative log-likelihood** as:
\begin{equation}
\mathcal{L}(U; \mathcal{W}) = -\sum_{i=1}^{N} \log p_{\mathcal{W}}(x_i \mid x_{\mathcal{D}_i}),
\label{eq:localNLL}
\end{equation}
where \( \mathcal{W} \) denotes the model parameters.
%
Let \( G \) be a directed graph where each edge is \( \mathcal{D}_i \to i \). We distinguish two cases:
%
\begin{enumerate}[noitemsep,nolistsep]
\item 
If \( G \) is a directed acyclic graph (DAG), then a topological ordering \( \tau \) exists such that \( \mathcal{D}_{\tau(i)} \subseteq \{\tau(1), \dots, \tau(i-1)\} \). In this case, the joint distribution $p_{\mathcal{W}}(x) = \prod_{i=1}^N p_{\mathcal{W}}(x_{\tau(i)} \mid x_{\mathcal{D}_{\tau(i)}})
$ is properly normalized and the local negative log-likelihood in \eqref{eq:localNLL} is equal to the true negative log-likelihood.
%
For example, during autoregressive training, each token is conditioned only on its predecessors, 
\begin{equation}
\mathcal{L}(U; \mathcal{W}) = -\sum_{i=1}^N \log p_{\mathcal{W}}(x_i \mid x_1, \dots, x_{i-1}) = -\sum_{i=1}^N \log p_{\mathcal{W}}(x_i \mid x_{<i}).
\label{eq:ARloss}
\end{equation}

\item If \( G \) contains directed cycles, the product $\prod_{i=1}^N p_{\mathcal{W}}(x_i \mid x_{\mathcal{D}_i})$ is not guaranteed to be normalized, and thus does not define a proper joint likelihood.
%
In this case, one can instead minimize the following pseudo log-likelihood,
\begin{equation}
\hat{\mathcal{L}}(U; \mathcal{W}) = -\sum_{i=1}^{N} \log p_{\mathcal{W}}(x_i \mid x_1, \dots, x_{i-1}, x_{i+1}, \dots, x_N) = -\sum_{i=1}^{N} \log p_{\mathcal{W}}(x_i \mid x_{-i}),
\label{eq:PLL}
\end{equation}
where \( x_{-i} \) denotes all tokens in the sequence except \( x_i \).
%
Following Besag's theorem \cite{besag1975statistical}, the maximum pseudo-likelihood estimator is consistent under standard regularity conditions, and \eqref{eq:PLL} remains compatible with stochastic-gradient
optimization even though it is not the true negative log-likelihood \cite{yang2019xlnet}.
\end{enumerate}





%% PROOF PROPOSITION GRADIENTS AND RANK-1 MATRIX
\subsection{Proof of Proposition \ref{prop-gradients-self-attention} and related remarks}
\label{supp-math-gradients-self-attention}
%
\input{ICML2025-rebuttal/math/proofs/proof-prop-gradient-self-attention}
%
In standard Transformer models, the bilinear form $\bm{W}_{qk}$ is not directly computed, and as such, it is not explicitly updated through gradient descent.
%
Nonetheless, $\bm{W}_{qk}$ is implicitly updated with a combination of the weight updates of $\bm{W}_q$ and $\bm{W}_k$ having the same form as in Proposition \ref{prop-gradients-self-attention}, see the following.
%
\input{ICML2025/math/statements/remark-self-attention-gradient-WqWk}



\subsection{Formal proofs of Theorem \ref{theo-gradient-directionality} and \ref{theo-gradients-symmetry}: the connection between objective functions, directionality, and symmetry}
\label{supp-math-directionality-symmetry}
%
In this section, we provide formal proofs of Theorems~\ref{theo-gradient-directionality} and~\ref{theo-gradients-symmetry}, which characterize the directional and symmetric structures that emerge during autoregressive and bidirectional training, respectively.
%
To this end, we introduce a series of intermediate Propositions and Lemmas, following these steps:

\begin{enumerate}[noitemsep,nolistsep]
    \item We begin by showing that when \( t^* \) is used as a context token to predict other tokens, all predicted tokens contribute to the \emph{column space} of \( \bm{W}_{qk} \), while only \( t^* \) contributes to the \emph{row space}.  
    Conversely, when \( t^* \) is the token being predicted, only \( t^* \) contributes to the column space, whereas all context tokens contribute to the row space (Section \ref{supp-math-prop-gradient-column-rows}).
%
    \item Next, we show that under reasonable assumptions about the statistical distribution of token embeddings, these structural properties imply that the expected norm of column updates exceeds that of row updates when \( t^* \) is used as context. In contrast, row updates dominate when \( t^* \) is the predicted token (Section \ref{supp-math-prop-gradient-asymmetric-growth-rows-columns}).
%
    \item We then show that in autoregressive training, the expected number of times a token \( t^* \) appears as context can differ from the expected number of times it is predicted, since these two quantities are affected differently by statistical correlations between tokens.
    %
    In contrast, under bidirectional training, the expected counts are always equal, regardless of token correlations (Section~\ref{supp-math-prop-counting-prediction-context}).
%
    \item Finally, we combine the above results to prove our main theorems:  
    (1) the weight updates to \( \bm{W}_{qk} \) induce \emph{column dominance} under autoregressive training (Section \ref{supp-math-lemma-tail-probabilities}-\ref{supp-math-theo-gradient-directionality}), and  
    (2) the weight updates to \( \bm{W}_{qk} \) induce \emph{symmetry} under bidirectional training (Section \ref{supp-math-theo-symmetry-gradients}).
\end{enumerate}
%

%
%% (1) PROPOSITION COLUMN AND ROWS
%
\subsubsection{Proof of Proposition \ref{prop-gradient-columns-rows}}
\label{supp-math-prop-gradient-column-rows}
%
First, we show that a token $t^*$ contributes differently to the updates of $\bm{W}_{qk}$ depending on whether it serves as context for predicting other tokens or is itself predicted.
%
\begin{proof}
%
Let $U = \{t_1, \dots, t_N\}$ be a sequence of tokens with the embedding of every $i$-th token be given by $\bm{x}_i \in \mathbb{R}^d$.
%
Proposition \ref{prop-gradients-self-attention} shows that the implicit weight update can be decomposed with two equivalent regrouping of the double summations, as follows,
%
\begin{equation}
\Delta \bm{W}_{qk} = \sum_{(i,j) \in U} \beta_{ij}\bm{x}_i\bm{x}_j^T 
= \sum_j \left(\sum_{i\in \mathcal{P}_j} \beta_{ij}\bm{x}_i\right)\bm{x}^\top_j
= \sum_i \bm{x}_i\left(\sum_{i\in \mathcal{C}_i} \beta_{ij}\bm{x}^\top_j\right)\,,
\end{equation}
%
where $\mathcal{P}_{i} \subset U$ is the set of tokens predicted by a given token $t_i$, while $\mathcal{C}_{j} \subset U$ is the set of tokens that predict a given token $t_j$.
%
%
We neglect any constant of proportionality - such as a learning rate - for simplicity, and we do not assume any specific structure on $\mathcal{P}_i$ and $\mathcal{C}_j$ (autoregressive training, bidirectional training, or others).
%
First, the contribution of $t^* \in U$ to the weight update when $t^*$ is used as context to predict a set of tokens $\mathcal{P}_{t^*} \subset U$ is
%%
\begin{equation}
\Delta \bm{W}_{qk} \bigg|_{t_j = t^*}
= \left(\sum_{i\in \mathcal{P}_{t^*}} \beta_{i*}\bm{x}_i\right)\bm{x}^\top_{t^*}\,.
\end{equation}
%
The associated weight update of the $k$-th column is then given by
%
\begin{equation}
    \Delta \bm{w}_{\cdot, k} \bigg|_{t_j = t^*} = 
    \left(\sum_{i\in \mathcal{P}_{t^*}} \beta_{i*}\bm{x}_i\right)[\bm{x}_{t^*}]_k
    =
    [\bm{x}_{t^*}]_k \left(\sum_{i \in \mathcal{P}_{t^*}}\beta_{i*}\bm{x}_i\right)\,,
\end{equation}
%
while the update of the $m$-th row is
%
\begin{equation}
    \Delta \bm{w}_{m, \cdot} \bigg|_{t_j = t^*} = \left(\sum_{i \in \mathcal{P}_{t^*}}\beta_{i*}[\bm{x}_i]_m\right)\bm{x}_{t^*} \,.
\end{equation}
%
A complementary argument can be made when predicting a given token $t^*$ from a set of tokens $\mathcal{C}_{t^*}$.
%
Indeed, the contribution to the weight update when $t^*$ is predicted by set of tokens $\mathcal{C}_{t^*} \subset U$ is given by 
%
\begin{equation}
\Delta \bm{W}_{qk} \bigg|_{t_i = t^*}
= \left(\sum_{j\in \mathcal{C}_{t^*}} \beta_{*j}\bm{x}_{t^*}\right)\bm{x}_j^\top\,.
\end{equation}
%
The associated weight update of the $k$-th column is then given by
%
\begin{equation}
    \Delta \bm{w}_{k, \cdot} \bigg|_{t_i = t^*} = \left(\sum_{j \in \mathcal{C}_{t^*}}\beta_{*j}[\bm{x}_j]_k\right)\bm{x}_{t^*} \,,
\end{equation}
%
while the update of the $m$-th row is
%
\begin{equation}
    \Delta \bm{w}_{m, \cdot} \bigg|_{t_i = t^*} =  \left(\sum_{j \in \mathcal{C}_{t^*}}\beta_{*j}\bm{x}_j^\top\right) [\bm{x}_{t^*}]_m
    = [\bm{x}_{t^*}]_m \left(\sum_{j \in \mathcal{C}_{t^*}}\beta_{*j}\bm{x}_j^\top\right)\,.
\end{equation}
%
Therefore, the weight update of each column (row) when a set of tokens predicts $t^*$ is equivalent to the weight update of each row (column) when $t^*$ is used to predict a set of tokens.
%
This concludes the proof.
%
\end{proof}
%

%





%%
%% (2) ASYMMETRIC GROWTH OF COLUMNS AND ROWS
\subsubsection{Asymmetric growth of rows and columns during weight update}
\label{supp-math-prop-gradient-asymmetric-growth-rows-columns}
%
Next, we demonstrate that under reasonable assumptions about the statistical distribution of token embeddings, the expected norm of column updates exceeds that of row updates when the token \( t^* \) is used as context to predict other tokens.
%
Conversely, when $t^*$ is being predicted by other tokens, the row updates become dominant.
%
To do so, we make the following assumptions:
%
\begin{itemize}[noitemsep,nolistsep]
%
    \item The token embeddings $\bm{x}_i$ are independent and identically distributed (i.i.d.) random vectors drawn from a probability distribution $\mathcal{D}$ with zero mean, $\mathbb{E}[\bm{x}_i] = \mathbf{0}$, and covariance matrix $\text{Cov}(\bm{x}_i) = \Sigma$.
    %
    This assumption holds at initialization for any Transformer model with learnable embeddings. 
    %
    \item The covariance matrix is not isotropic, i.e., $\Sigma \neq \sigma^2 \mathbb{I}$,.
    %
    More specifically, we posit that there is partial alignment between the embeddings $\bm{x}_i$ due to the semantic and predictive relationships between tokens, which typically emerge during training.
%
\end{itemize}
%
Similar to Proposition \ref{prop-gradient-columns-rows}, the scenarios where $t^*$ is used to predict other tokens and where $t^*$ is being predicted by other tokens are complementary. 
%
In the following Proposition, we focus solely on the case where $t^*$ serves as context to predict a set of tokens. 
%
A formal derivation for the opposite case where $t^*$ is predicted by other tokens is provided in a subsequent Corollary.

\begin{proposition}
\label{prop-gradient-asymmetric-growth-rows-columns}
%
(\textbf{Asymmetric growth of columns and rows for context}).
%
Let $U = \{t_1, \dots, t_N\}$ a sequence of tokens, and let $t^*$ be a token representing the context of every token $t_i \in U$.
%
Let $\{\bm{x}_1 \dots, \bm{x}_N\}$ be the token embedding associated with $U$ such that each $\bm{x}_i \sim \mathcal{D}$ is a i.i.d. random vector drawn from a probability distribution $\mathcal{D}$ with zero mean and non-isotropic covariance $\text{Cov}(\bm{x}_i)  = \Sigma$.
%
Let $\Delta \bm{W}_{qk}$ be the weight update from Proposition \ref{prop-gradients-self-attention}.
%
Then the squared norm of the $m$-th row and $k$-th column of $\bm{W}_{qk}$ satisfies
%
\begin{equation}
\frac{\mathbb{E}_{\mathcal{D}}\left[ ||\Delta \bm{w}_{\cdot, k}||^2 \right]}{\mathbb{E}_{\mathcal{D}}\left[ ||\Delta \bm{w}_{m, \cdot}||^2 \right]} > 1 \quad \forall k \in \{1,\dots,d\}, \forall m \,\,\text{s.t.} \,\, \Sigma_{m,m} < \frac{\text{Tr}(\Sigma)}{d}
\end{equation}
%
\end{proposition}
\input{ICML2025-rebuttal/math/proofs/proof-prop-gradient-asymmetric-growth-rows-columns}
%

%
Again, a complementary argument can be made about the asymmetric weight update when a set of tokens $U = \{t_1, \dots, t_N\}$ is used to predict a given token $t^*$. We formalize this in the following Corollary.
%
\input{ICML2025-rebuttal/math/statements/corollary-gradient-asymmetric-growth-rows-columns}
\begin{proof}
%
Let $U = [t_0, \dots, t_N]$ a sequence of tokens, and let $t^*$ be predicted by every token $t_i \in U$. 
%
Let $\{\bm{x}_i\}$ be the set of token embeddings associated with $U$ and let $\bm{y}$ be the token embedding of $t^*$.
%
It follows from Proposition \ref{prop-gradient-columns-rows} that the squared norm of the weight update of the $k$-th column is given by 
%
\begin{equation}
\Delta \bm{w}_{\cdot, k} =  \sum_{j=1}^N \beta_{*j}[\bm{x}_j]_k\bm{y} 
%
\,\,\Rightarrow \,\,
%
||\Delta \bm{w}_{\cdot, k}||^2 = \big(\sum_{i=1}^N \beta_{*j}[\bm{x}_j]_k\big)^2 ||  \bm{y}||^2 \,,
\end{equation}
%
while the weight update of the $m$-th row follows
%
\begin{equation}
\Delta \bm{w}_{m, \cdot} = \sum_{j=1} ^N \beta_{*j} \bm{x}_j[\bm{y}]_m 
\,\,\Rightarrow \,\,
||\Delta \bm{w}_{m, \cdot}||^2 = [\bm{y}]^2_m || \sum_{i=1}^N \beta_{*j} \bm{x}_j||^2 \,.
\end{equation}
%
Therefore, following the same arguments as in Proposition \ref{prop-gradient-asymmetric-growth-rows-columns}, the ratio of the expected value of the square norms is given by
%
\begin{equation}
\frac{\mathbb{E}_{\mathcal{D}}\left[ ||\Delta \bm{w}_{\cdot, k}||^2 \right]}{\mathbb{E}_{\mathcal{D}}\left[ ||\Delta \bm{w}_{m, \cdot}||^2 \right]} = \frac{ \Sigma_{k,k}\text{Tr}(\Gamma)}{ \Gamma_{m,m}\text{Tr}(\Sigma)} = d\frac{\Sigma_{k,k}}{\text{Tr}(\Sigma)}  < \frac{d\,\Sigma_{k,k}}{d\,\Sigma_{k,k}} = 1 \quad \forall k \,\,\text{s.t.}\,\,\Sigma_{k,k} < \frac{\text{Tr}(\Sigma)}{d} \,,
\end{equation}
%
thus concluding the proof.
%
%
\end{proof}






%
%% (3) PROPOSITION COUNTING + PROOF
\subsubsection{Expected contribution for context and prediction and related remarks}
\label{supp-math-prop-counting-prediction-context}
%%
%%
%% PROVIDE ASSUMPTIONS ON MASKING AND BIDIRECTIONAL TRAINING
%
Additionally, we show that in autoregressive training, the expected number of tokens predicted by a given token \( t^* \) can differ from the number of tokens that predict \( t^* \), with this ratio influenced by token correlations, while in bidirectional training, the two quantities are always equal regardless of such correlations, yielding a ratio of 1.
%
We formalize this in the following proposition and illustrate, in the subsequent remark, an example from autoregressive training where the expected ratio exceeds 1.
%
\begin{proposition}
\label{prop-counting-prediction-context}
%
(\textbf{Ratio of expected counts during autoregressive and bidirectional training}).
%
Let $\mathcal{V} = \{t_0, \dots, t_V\}$ be a set of tokens.
%
Let $\mathcal{U}$ be the sample space of all possible sequences of $N$ tokens, and let the sequence $U = \{t_1, \dots, t_N\} \in\mathcal{U}$ be a random variable with probability distribution $P(U)$ defined over $\mathcal{U}$.
%
Let $\Pr[t_j = t^*]$ be the probability that the token at index $j$ in a sequence $U \in \mathcal{U}$ is given by $t^* \in \mathcal{V}$.
%
Let $\mathbb{E}_{P}[\mu_c(t^*)]$ be the expected number of tokens that are predicted by a given token $t^*$, and $\mathbb{E}_{P}[\mu_p(t^*)]$ the expected number of tokens that predict a given token $t^*$.
%
For autoregressive training, the ratio
%
\begin{equation}
    \frac{\mathbb{E}_{P}[\mu_c(t^*, U)]}{\mathbb{E}_{P}[\mu_p(t^*, U)]} = \frac{\sum_{k=1}^N (N-k)\Pr[t_k = t^*]}{\sum_{k=1}^N(k-1)\Pr[t_k = t^*]} \,,
\end{equation}
%
depends on $\Pr[t_k = t^*] \,\,\forall t^* \in \mathcal{V}, \forall k \in \{1,\dots. N\}$, while for bidirectional training the same ratio is given by 
%
\begin{equation}
     \frac{\mathbb{E}_{P}[\mu_c(t^*, U)]}{\mathbb{E}_{P}[\mu_p(t^*, U)]} = 1 \,.
\end{equation}
%    
\end{proposition}

% \begin{proposition}[Ratio of expected counts]
% \label{prop:count_ratio}
% %
% Fix a vocabulary $\mathcal V=\{t_0,\dots,t_V\}$ and an integer
% sequence length $N\ge 2$.
% Let
% \(
%   (\Omega,\mathscr F,\Pr)
% \)
% be a probability space that carries  

% * random tokens
%   \(
%     T_1,\dots,T_N:\Omega\to\mathcal V,
%   \)
%   collected in the random sequence
%   \(U=(T_1,\dots,T_N)\in\mathcal V^{N}\), and  

% * (for the bidirectional‐training analysis) i.i.d.\
%   Bernoulli masking variables
%   \(
%     M_1,\dots,M_N\sim\mathrm{Bernoulli}(\rho),
%     \quad 0<\rho<1,
%   \)
%   which are independent of each other and of the tokens.

% For a fixed token
% \(t^*\in\mathcal V\) define the random variables
% \begin{align}
%   \mu_c(t^*,U)
%     &=\sum_{k=1}^{N}(N-k)\,
%        \mathbf 1_{\{T_k=t^*\}}, &\text{(tokens \emph{predicted by} }t^*),
%   \\
%   \mu_p(t^*,U)
%     &=\sum_{k=1}^{N}(k-1)\,
%        \mathbf 1_{\{T_k=t^*\}}, &\text{(tokens that \emph{predict} }t^*).
% \end{align}
% Write
% \(
%   E[\cdot]=\E_{U,M_1{:}N}[\cdot]
% \)
% for expectation with respect to all randomness.

% \begin{enumerate}[label=(\roman*),topsep=4pt,itemsep=2pt]
% \item \textbf{Autoregressive training.}
%       Then
%       \begin{equation}\label{eq:ratio_ar}
%         \frac{E[\mu_c(t^*)]}{E[\mu_p(t^*)]}
%         \;=\;
%         \frac{\displaystyle\sum_{k=1}^{N}(N-k)\Pr[T_k=t^*]}
%              {\displaystyle\sum_{k=1}^{N}(k-1)\Pr[T_k=t^*]}\;.
%       \end{equation}

% \item \textbf{Bidirectional (masked-LM) training.}
%       Assume the model predicts every \emph{masked} token using
%       every \emph{unmasked} token.
%       Under the independence assumptions stated above,
%       \[
%         \frac{E[\mu_c(t^*)]}{E[\mu_p(t^*)]}\;=\;1 .
%       \]
% \end{enumerate}
% \end{proposition}

% \begin{proof}
% \textbf{(i) Autoregressive case.}
% Because $\mathbf 1_{\{T_k=t^*\}}$ is
% $\{0,1\}$‐valued, linearity of expectation yields
% \[
%   E[\mu_c(t^*)]
%   =\sum_{k=1}^{N}(N-k)E[\mathbf 1_{\{T_k=t^*\}}]
%   =\sum_{k=1}^{N}(N-k)\Pr[T_k=t^*].
% \]
% The same calculation with the weights $(k-1)$ gives
% \(
%   \E[\mu_p(t^*)]=\sum_{k=1}^{N}(k-1)\Pr[T_k=t^*].
% \)
% Taking the ratio gives \eqref{eq:ratio_ar}.

% \medskip
% \noindent
% \textbf{(ii) Bidirectional case.}
% Fix a position $k$.
% If $T_k=t^*$ and $M_k=0$ (token visible),
% it contributes
% \(
%   \sum_{k'\ne k}\mathbf 1_{\{M_{k'}=1\}}
% \)
% to $\mu_c(t^*)$.
% If $T_k=t^*$ and $M_k=1$ (token masked),
% it contributes
% \(
%   \sum_{k'\ne k}\mathbf 1_{\{M_{k'}=0\}}
% \)
% to $\mu_p(t^*)$.
% Taking expectations and using independence of
% \(\{M_i\}\) together with
% \(E[\mathbf 1_{\{M_i=1\}}]=\rho\),
% \(E[\mathbf 1_{\{M_i=0\}}]=1-\rho\),
% we obtain
% \[
%   E[\mu_c(t^*)]
%   =\sum_{k=1}^{N}\Pr[T_k=t^*]\,(1-\rho)\,(N-1)\rho,
%   \qquad
%   E[\mu_p(t^*)]
%   =\sum_{k=1}^{N}\Pr[T_k=t^*]\,(  \rho)\,(N-1)(1-\rho).
% \]
% The multiplicative factors are identical, hence
% \(
%   E[\mu_c(t^*)]=E[\mu_p(t^*)]
% \)
% and the stated ratio equals~$1$.
% \end{proof}

% \begin{remark}[Effect of token–position statistics]\label{rem:iid_vs_corr}
% \hfill
% \begin{enumerate}[label=(\alph*),leftmargin=*,topsep=2pt,itemsep=2pt]
% \item \emph{i.i.d.\ tokens.}
%       If the tokens are independent and identically distributed,
%       \(\Pr[T_k=t^*]=\pi_{t^*}\) for all $k$.
%       Then the two sums in \eqref{eq:ratio_ar} coincide,
%       giving
%       \(
%         E[\mu_c(t^*)]=E[\mu_p(t^*)]=
%         \pi_{t^*}\tfrac{N(N-1)}2
%       \)
%       and the ratio is~1.  Hence no directional bias appears.

% \item \emph{Correlated tokens.}
%       For a general joint law
%       \(\Pr[T_1{:}T_N]\),
%       \eqref{eq:ratio_ar} remains valid but cannot be simplified.
%       Writing the joint law via the reverse chain rule
%       \(
%         \Pr[T_1{:}T_N]
%         =\prod_{i=1}^{N}\Pr[T_i\mid T_{i+1{:}N}],
%       \)
%       one obtains
%       \[
%         \Pr[T_k=t^*]
%         =\!\!\!\sum_{\substack{t_{-k}\in\mathcal V^{N-1}}}
%           \Bigl(
%             \Pr[T_k=t^*\mid T_{k+1{:}N}=t_{k+1{:}N}]
%             \!\!\!\prod_{i\ne k}
%             \Pr[T_i=t_i\mid T_{i+1{:}N}=t_{i+1{:}N}]
%           \Bigr),
%       \]
%       which may depend strongly on $k$.
%       If $t^*$ is more likely to occur at early positions
%       than late ones, the numerator of \eqref{eq:ratio_ar}
%       (weighted by $N-k$) becomes larger than the denominator,
%       giving a ratio \(>1\) and
%       inducing directionality even under autoregressive training.
% \end{enumerate}
% \end{remark}

\begin{proof}
%
Let $\mathcal{V} = \{t_0, \dots, t_V\}$ be a set of tokens.
%
Let \( \mathcal{U} = \mathcal{V}^N \) denote the sample space of all possible sequences of \( N \) tokens from the vocabulary \( \mathcal{V} \). 
%
Let \( \mathcal{F} = 2^{\mathcal{U}} \) be the \(\sigma\)-algebra given by the power set of \( \mathcal{U} \), and let \( P \) be the probability distribution over \( \mathcal{U} \) that defines the distribution of the random variable sequence \( U \). 
%
This defines the probability space \( (\mathcal{U}, \mathcal{F}, P) \).
%
For each position $k \in \{1, \dots, N\}$ we let the random variable $T_k$ be $T_k: \mathcal{U} \rightarrow \mathcal{V}$ such that $T_k(U) \equiv t_k$.
%
Let $E = \{U \in \mathcal{U}: t_k = t^*\}$ the event of all sequences where the $k$-th token is a fixed token $t^*$. 
%
The indicator function $\mathds{1}\{E\}(U)$ is a random variable defined as
%
\begin{equation}
\mathds{1}\{E\}(U) = \mathds{1}\{t_k = t^*\}(U) =
\begin{cases}
    1 \,\,\text{if} \,\,t_k = t^*\\
    0 \,\, \text{otherwise} \,,
\end{cases}
\end{equation}
%
and as such its expected value over $P$ is
%
\begin{equation}
\mathbb{E}_P\left[\mathds{1}\{t_k = t^*\}(U)\right] = \sum_{U \in \,\mathcal{U}} \mathds{1}\{t_k = t^*\}(U)\Pr[U] = \sum_{U \in \mathcal{U} \,:\, t_k = t^*} \Pr[U] = \Pr[t_k = t^*] \,.
\end{equation}
%  
Let $\mu_c(t^*, U)$ be the random variable quantifying the number of tokens predicted by $t^*$, while $\mu_p(t^*, U)$ is the random variable quantifying the number of tokens predicted by $t^*$.
%
We analyzed the case of autoregressive and bidirectional training separately.
%

During autoregressive training, each time the token $t^*$ appears at position $k$ it is used as context to predict $N-k$ tokens, and it is predicted by $k-1$ tokens.
%
It follows that $\mu_c(t^*, U)$ is given by 
%
\begin{equation}
\mu_c(t^*, U) = \sum_{l=2}^N\sum_{k=1}^{l-1} \mathds{1}\{t_k = t^*\} = \sum_{k=1}^N (N-k)\mathds{1}\{t_k = t^*\}\,,
\end{equation}
%
while $\mu_p(t^*, U)$ is given by 
%
\begin{equation}
\mu_p(t^*, U) = \sum_{l=1}^{N-1}\sum_{k=1}^{l-1} \mathds{1}\{t_l = t^*\} = \sum_{k=1}^N (k-1)\mathds{1}\{t_k = t^*\}\,.
\end{equation}
%
Therefore, the expected value of $\mu_c(t^*, U)$ over $P$ is
%
\begin{equation}
\mathbb{E}_P[\mu_c(t^*, U)] = \mathbb{E}_P\left[\sum_{k=1}^N (N-k)\mathds{1}\{t_k = t^*\}\right] =   \sum_{k=1}^N(N-k)\Pr[t_k = t^*]\,,
\end{equation}
%
the expected value of $\mu_p(t^*)$ is
%
\begin{equation}
\mathbb{E}_P[\mu_p(t^*, U)] = \mathbb{E}_P\left[\sum_{k=1}^N (k-1)\mathds{1}\{t_k = t^*\}\right] =   \sum_{k=1}^N(k-1)\Pr[t_k = t^*]\,,
\end{equation}
%
and their ratio is given by
%
\begin{equation}
\label{suppeq:ratio-autoregressive}
    \frac{\mathbb{E}_P[\mu_c(t^*, U)]}{\mathbb{E}_P[\mu_p(t^*, U)]} = \frac{\sum_{k=1}^N (N-k)\Pr[t_k = t^*]}{\sum_{k=1}^N(k-1)\Pr[t_k = t^*]} \,.
\end{equation}
%

During bidirectional training, each time the token $t^*$ is masked at position $k$, it is predicted by $N-1$ tokens.
%
We assume that the token $t^*$ significantly contributes to the context of a masked token only if it is itself not masked (see Appendix \ref{supp-math-self-attention-derivation-gradients}).
%
Additionally, we assume that the probability $\rho$ of masking any given position $k$ is the same for all positions and does not depend on the specific token $t_k$ at that position.
%
As such, we define masking as an independent Bernoulli variable $m_k: \Sigma_k \rightarrow \{0.1\} \sim P_m$ such that $\text{Pr}[m_k = 1] = \rho$, independently of the token and sequence spaces.
%
It follows that $\mu_c(t^*, U, m)$ is given by 
%
\begin{equation}
\mu_c(t^*, U, m) = \sum_{k=1}^N\mathds{1}\{t_k = t^*\} \mathds{1}\{m_k = 0\}\sum_{k' \neq k} \mathds{1}\{m_{k'} = 1\}  \,,
\end{equation}
%
while $\mu_p(t^*, U, m)$ is given by 
%
\begin{equation}
\mu_p(t^*, U, m) = \sum_{k=1}^N\mathds{1}\{t_k = t^*\} \mathds{1}\{m_{k} = 1\}\sum_{k' \neq k} \mathds{1}\{m_{k'} = 0\}\,.
\end{equation}
%
Therefore, by taking the expectation with respect to both distributions \( P \) and \( P_m \), their respective expected values are given by
%
\begin{equation}
\begin{split}
\mathbb{E}_{P,P_m}[\mu_c(t^*, U, m)] 
& =  \mathbb{E}_{P,P_m} \left[\sum_{k=1}^N\mathds{1}\{t_k = t^*\} \mathds{1}\{m_{k} = 0\}\sum_{k' \neq k} \mathds{1}\{m_{k'} = 1\} \right] \\
& = \sum_{k=1}^N \mathbb{E}_{P,P_m} \left[\mathds{1}\{t_k = t^*\} \mathds{1}\{m_{k} = 0\}\sum_{k' \neq k} \mathds{1}\{m_{k'} = 1\} \right]\\
& = \sum_{k=1}^N \mathbb{E}_{P,P_m} \left[\mathds{1}\{t_k = t^*\} \mathds{1}\{m_{k} = 0\}\right]\sum_{k' \neq k} \left[\mathds{1}\{m_{k'} = 1\} \right]\\
& = \sum_{k=1}^N \text{Pr}(t_k = t^*) \text{Pr}(m_{k} = 0)\sum_{k' \neq k} \text{Pr}(m_{k'} = 1) \\
& = N\rho(N-1)(1 - \rho) \sum_{k=1}^N \text{Pr}(t_k = t^*)\,,
\end{split}
\end{equation}
%
and
%
\begin{equation}
\begin{split}
\mathbb{E}_{P,P_m}[\mu_p(t^*, U, m)] 
& =  \mathbb{E}_{P,P_m} \left[\sum_{k=1}^N\mathds{1}\{t_k = t^*\} \mathds{1}\{m_{k} = 1\}\sum_{k' \neq k} \mathds{1}\{m_{k'} = 0\} \right] \\
& = N\rho(N-1)(1 - \rho)\sum_{k=1}^N \text{Pr}(t_k = t^*)  \,,
\end{split}
\end{equation}
%
and their ratio is
\begin{equation}
     \frac{\mathbb{E}_{P}[\mu_c(t^*, U)]}{\mathbb{E}_{P}[\mu_p(t^*, U)]} = 1 \,,
\end{equation}
%
where we omit the dependence on the independent random variable $m$ for simplicity, thus concluding the proof.
%
\end{proof}
%

The following remark highlights how, during autoregressive training, the ratio of expected token counts can vary depending on the correlations and conditional dependencies among the tokens.
%
\begin{remark}
\label{remark:ratio-autoregressive}
%
We separately analyze the following two cases:
%
\begin{itemize}[noitemsep,nolistsep]
%
    \item \emph{No statistical correlation between tokens.}
    %
    Assume that the tokens are statistically independent and identically distributed (i.i.d.) across positions, i.e., the probability of each token is the same at every position and does not depend on surrounding tokens.
    %
    Under this assumption, during autoregressive training, the expected number of times token \( t^* \) is used to predict future tokens is given by,
    %
    \begin{equation}
    \mathbb{E}_P[\mu_c(t^*, U)] = \sum_{k=1}^N (N - k)\Pr[t_k = t^*] = \Pr[t^*] \sum_{k=1}^N (N - k) = \Pr[t^*] \frac{N(N-1)}{2} \,.
    \end{equation}
    %
    Similarly, the expected number of times token \( t^* \) is predicted by preceding tokens is,
    \begin{equation}
    \mathbb{E}_P[\mu_p(t^*, U)] = \sum_{k=1}^N (k - 1)\Pr[t_k = t^*] = \Pr[t^*] \sum_{k=1}^N (k - 1) = \Pr[t^*] \frac{N(N-1)}{2} \,.
    \end{equation}
    %
    Therefore, when tokens are independent and identically distributed, the expected number of tokens predicted \emph{by} \( t^* \) and the expected number of tokens used to predict \( t^* \) are equal. In this case, there is no difference between autoregressive and bidirectional training.
    %
    \item \emph{Statistical correlation between tokens.} 
    %
    Let us assume that the probability distribution \( P \) encodes statistical dependencies between tokens.
    %
    Define \( \mu(t^*) \) as the expected position index of token \( t^* \),
    \begin{equation}
        \mu(t^*) = \frac{\sum_{k=1}^N k \Pr[t_k = t^*]}{\sum_{k=1}^N \Pr[t_k = t^*]} \,.
    \end{equation}
    %
    Using this, the ratio of expected counts under autoregressive training becomes,
    \begin{equation}
    \frac{\mathbb{E}_P[\mu_c(t^*, U)]}{\mathbb{E}_P[\mu_p(t^*, U)]} = \frac{N - \mu(t^*)}{\mu(t^*) - 1} \,.
    \end{equation}
    %
    This expression shows that if \( t^* \) tends to appear earlier in the sequence, that is, if \( \mu(t^*) < (N+1)/2 \) (left of center), then the ratio is greater than 1. 
    %
    Conversely, if \( t^* \) appears later in the sequence (right of center), the ratio is less than 1.
    %
    More formally, suppose there exists \( \delta > 0 \) such that
    \begin{equation}
        \mu(t^*) \le \frac{N+1}{2} - \delta \,.
    \end{equation}
    % 
    Then the ratio can be rewritten as,
    \begin{equation}
        \frac{\mathbb{E}_P[\mu_c(t^*, U)]}{\mathbb{E}_P[\mu_p(t^*), U]} - 1
        = \frac{N + 1 - 2\mu(t^*)}{\mu(t^*) - 1} \,,
    \end{equation}
    where the numerator satisfies \( N + 1 - 2\mu(t^*) \ge 2\delta \), and the denominator is positive for any token that can appear at position 1 (i.e., \( \mu(t^*) > 1 \)).
    %
    Hence, the ratio is bounded below by,
    \begin{equation}
        \frac{\mathbb{E}_P[\mu_c(t^*, U)]}{\mathbb{E}_P[\mu_p(t^*, U)]} \ge 1 + \varepsilon(t^*) \,, \quad
        \text{with} \quad \varepsilon(t^*) = \frac{2\delta}{\mu(t^*) - 1} \,.
    \end{equation}
%
\end{itemize}
%
\end{remark}
%

% %
% We provide some examples where tokens that can tend to appear earlier in the sequence.
% %
% As a first example, consider a position-biased mixture distribution \( P \), where tokens are independent across positions, but the probability of observing a specific token \( t^* \) depends on the position \( k \). Specifically, \( \Pr[t_k = t^*] = \alpha q_k \), where \( \alpha \in (0,1] \) and \( \{q_k\} \) is a non-increasing sequence in \( k \).
% %
% For instance, \( q_k \) may follow a linear decay \( q_k \propto (N + 1 - k) \), or an exponential decay \( q_k \propto e^{-\lambda(k - 1)} \), such that \( q_1 \ge q_2 \ge \cdots \ge q_N \). 
% %
% In such cases, any token whose probability decreases with position \( k \) will exhibit an expected count ratio greater than 1.
% %
% Next, consider a distribution \( P \) defined by a Zipf-tail model governed by a two-state hidden Markov chain. Assume that in state 0, the token \( t^* \) is emitted with a small probability \( \theta \), while in state 1, it is emitted with a higher probability \( \Theta > \theta \), and the chain remains in state 1 with probability \( \varphi \).
% %
% Additionally, suppose that bursts are more likely to begin near the start of the sequence. For example, the initial state is set to 1 with probability \( p_{\mathrm{init}} \propto k^{-\beta} \), decaying with \( k \) according to a Zipf law.


%% (5.0) LONG-TAIL DISTRIBUTION AND HIGHER NUMBER OF OUTLIERS
\subsubsection{Tail probabilities of distributions with equal mean and different variances}
\label{supp-math-lemma-tail-probabilities}
%
Before proving column dominance in the next Theorem, we first establish that a probability distribution with higher variance is more likely to produce large values.  
%
Specifically, when two distributions share the same mean but differ in variance, the one with higher variance has a greater probability of generating samples that exceed a given threshold.  
%
We formalize this in the following lemma.
%
%
\begin{lemma}
\label{lemma-same_mean_different_variance_bound}
Let \( a \) and \( b \) be two probability distributions with the same mean \( \mu \), but different variances such that \( \sigma_a^2 > \sigma_b^2 \). 
%
Then, for all values \( z > \sqrt{\sigma_a \sigma_b} - \mu \), the probability that a random sample from distribution \( a \) exceeds \( z \) is greater than or equal to the probability that a sample from distribution \( b \) exceeds \( z \), that is,
%
\begin{equation}
    \Pr[X_a > z] \geq \Pr[X_b > z] \quad \text{for all } z > \sqrt{\sigma_a \sigma_b} - \mu \,.
\end{equation}
%
\end{lemma}
%
\input{ICML2025-rebuttal/math/proofs/proof_lemma-probability_same_mean_different_variance}
%

%% (5) THEOREM DIRECTIONALITY
\subsubsection{Proof of Theorem \ref{theo-gradient-directionality}}
\label{supp-math-theo-gradient-directionality}
%
Here, we show that the weight updates of $\bm{W}_{qk}$ induce column dominance during autoregressive training, leading to the emergence of directionality.
%
To do so, we make the following assumptions:
%
\begin{itemize}[noitemsep,nolistsep]
%
    \item We adopt the same assumptions stated in Propositions~\ref{prop-gradient-asymmetric-growth-rows-columns} and~\ref{prop-counting-prediction-context}.
    %
    \item At initialization, the entries of \( \bm{W}_{qk} \) are drawn from a probability distribution \( \mathcal{P} \) with finite mean \( \mu \) and variance \( \sigma^2 \).  
    This is satisfied by any standard machine learning initialization scheme. This second assumption implies that each $k$-th column $\bm{w}_{\cdot,k}$ and $m$-th row $\bm{w}_{m, \cdot}$ have the same mean $\mu$ and variance $\sigma^2/\sqrt{n}$.
    %
    Furthermore, the squared norm of both rows and columns has equal mean $n(\sigma^2 + \mu^2)$ and variance $n(2\sigma^4 + 4\mu^2)$.
%
\end{itemize}
%
\begin{proof}
%
Let $\{U_1, \dots, U_D\}$ be a dataset of sequences of length $N$ drawn i.i.d. from a distribution $P$.
%
Let \( k_u \) denote the position of token \( t^* \) in sequence \( U_u \), and let superscript \( u \) refer to tokens in the \( u \)-th sequence.
%
Following Proposition \ref{prop-gradient-columns-rows} and \ref{prop-gradient-asymmetric-growth-rows-columns}, we obtain 
%
\begin{equation}
\Delta \bm{w}_{\cdot, k}  =  \sum_{u=1}^D \left[[\bm{y}]_k \sum_{i=k_u+1}^N \beta^u_{i*}\bm{x}^u_i + \sum_{j=1}^{k_u-1} \beta^u_{*j}[\bm{x}^u_j]_k\bm{y}
\right] \,,
\end{equation}
%
while the squared norm of the weight update of the $m$-th row by
%
\begin{equation}
\Delta \bm{w}_{m, \cdot}  =  \sum_{u=1}^D \left[ \sum_{i=k_u+1}^N \beta^u_{i*}[\bm{x}^u_i]_m\bm{y} + [\bm{y}]_m\sum_{j=1}^{k_u-1} \beta^u_{*j}\bm{x}^u_j
\right]\,,
\end{equation}
%
where in both equations the superscript $u$ indicates tokens belonging to the $u$-th sequence in the dataset, and $k_u$ indicates the position index of $t^*$.
%
It follows again from Proposition \ref{prop-gradient-asymmetric-growth-rows-columns} that the expected value of the squared norm of these updates is
%
\begin{equation}
\label{suppeq:ratio-squared-norm-theorem-directionality}
\frac{\mathbb{E}_{\mathcal{D}}\left[ ||\Delta \bm{w}_{\cdot, k}||^2 \right]}{\mathbb{E}_{\mathcal{D}}\left[ ||\Delta \bm{w}_{m, \cdot}||^2 \right]} \approx \frac{\Gamma_{kk}\text{Tr}(\Sigma)\big(\sum_{u=1}^D\sum_{i=k_u+1}^N \beta^u_{i*}\big) + \Sigma_{kk}\text{Tr}(\Gamma)\big(\sum_{u=1}^D\sum_{j=1}^{k_u-1} \beta^u_{*j}\big)}{\Sigma_{mm}\text{Tr}(\Gamma)\big(\sum_{u=1}^D\sum_{i=k_u+1}^N \beta^u_{i*}\big) + \Gamma_{mm}\text{Tr}(\Sigma)\big(\sum_{u=1}^D\sum_{j=1}^{k_u-1} \beta^u_{*j}\big)} \,.
\end{equation}
%
Following Proposition \ref{prop-counting-prediction-context}, let $\mu_c(t^*, U)$ be the random variable quantifying the number of tokens predicted by $t^*$, while $\mu_p(t^*, U)$ is the random variable quantifying the number of tokens predicted by $t^*$.
%
When assuming autoregressive training, it follows that the empirical average of $\mu_c(t^*, U)$ and $\mu_p(t^*, U)$ over the dataset $\{U_1, \dots, U_D\}$ approaches,
%
\begin{equation}
    \frac{\mathbb{E}_P[\mu_c(t^*, U)]}{\mathbb{E}_P[\mu_p(t^*, U)]} = \frac{\sum_{j=1}^N (N-j)\Pr[t_j = t^*]}{\sum_{j=1}^N(j-1)\Pr[t_j = t^*]} \,.
\end{equation}
%
Following Remark \ref{remark:ratio-autoregressive}, assume that the expected position \( \mu(t^*) \) of \( t^* \) satisfies \( \mu(t^*) < (N+1)/2 - \delta \) for some \( \delta > 0 \), then
%
\begin{equation}
\frac{\mathbb{E}_P[\mu_c(t^*, U)]}{\mathbb{E}_P[\mu_p(t^*, U)]} > 1 + \epsilon(t^*)\,,
\end{equation}
%
for some \( \epsilon(t^*) > 0 \).
%
Hence, across the dataset, \( t^* \) appears more often as a context token than as a predicted token. 
%
This implies \( \mathbb{E}_P[k_u] < (N+1)/2 \), so that the first terms in the numerator and denominator of Equation~\eqref{suppeq:ratio-squared-norm-theorem-directionality} dominate, leading to, 
%
\begin{equation}
\frac{\mathbb{E}\left[ ||\Delta \bm{w}_{\cdot, k}||^2 \right]}{\mathbb{E}\left[ ||\Delta \bm{w}_{m, \cdot}||^2 \right]} > 1 \,,
\end{equation}
%
that is, the net increase in column norms exceeds that of row norms.
%
Assume \( \bm{W}_{qk} \) is initialized with i.i.d. entries from a distribution with mean \( \mu \) and variance \( \sigma^2 \), then at the beginning of training
$\mathrm{Var}(\|\bm{w}_{\cdot, k}\|) = \mathrm{Var}(\|\bm{w}_{m, \cdot}\|)
$, while after training we have
%
$
\mathrm{Var}(\|\bm{w}_{\cdot, k}\|) > \mathrm{Var}(\|\bm{w}_{m, \cdot}\|) \,\,\forall k, m.
$
Applying Lemma~\ref{lemma-same_mean_different_variance_bound}, it follows that for all \( w > \gamma \), where \( \gamma = \sqrt{ \mathrm{Var}(\|\bm{w}_{\cdot, k}\|) \cdot \mathrm{Var}(\|\bm{w}_{m, \cdot}\|)} - \mu \),
\begin{equation}
    \text{Pr}[||\bm{w}_{\cdot, k}|| > w] > \text{Pr}[|| \bm{w}_{m, \cdot}|| > w] \,\,\ \forall \,w > \gamma \,,
\end{equation}
% 
thus concluding the proof.
%
\end{proof}
%







%% THEOREM SYMMETRY
\subsubsection{Proof of Theorem \ref{theo-gradients-symmetry} and related remarks}
\label{supp-math-theo-symmetry-gradients}
%
Finally, we prove that the weight updates of $\bm{W}_{qk}$ induce symmetry during bidirectional training.
%
Importantly, the column dominance is present only during autoregressive training.
%
Indeed, it follows from Proposition \ref{prop-counting-prediction-context} that, during bidirectional training, the net increase of the norm of the columns is equal to the net increase of the norm of the rows.
%
We formalize this in the following Corollary.
%
\begin{corollary}
%
(\textbf{Bidirectional training does not induce directionality})
%
Let \( \mathcal{V} = \{t_0, \dots, t_V\} \) be a vocabulary of tokens, and let \( \mathcal{U} \subset \mathcal{V}^N \) denote the sample space of all sequences of length \( N \). Let \( U = \{t_1, \dots, t_N\} \in \mathcal{U} \) be a random variable with \( U \sim P(U) \). Define \( \Pr[t_j = t^*] \) as the marginal probability that the token at position \( j \) equals \( t^* \in \mathcal{V} \).
%
Let \( \{\bm{x}_1, \dots, \bm{x}_N\} \) be the token embeddings corresponding to the elements of \( U \), where each embedding \( \bm{x}_i \sim \mathcal{D} \) is drawn i.i.d. from a distribution \( \mathcal{D} \) with zero mean and covariance matrix $\text{Cov}(\bm{x}_i)  = \Sigma$..
%
Let \( \bm{W}_{qk} \) be query-key matrix of a self-attention mechanism, and let \( \Delta \bm{W}_{qk} \) denote its gradient update as defined in Proposition~\ref{prop-gradients-self-attention}, computed under an autoregressive objective as in Definition~\ref{def-objective-functions}.
%
Then, the variance of the norm of the rows and the columns at the end of training is equal,
%
%
\begin{equation}
    \text{Var}\left(||\bm{w}_{m, \cdot}||\right) =\text{Var}\left(||\bm{w}_{ \cdot, k}||\right) \,\, \forall k,m \,.
\end{equation}
%
\end{corollary}
\begin{proof}
%
%
Let $\{U_1, \dots, U_D\}$ be a dataset of sequences of length $N$ drawn i.i.d. from a distribution $P$.
%
Let \( k_u \) denote the position of token \( t^* \) in sequence \( U_u \), and let superscript \( u \) refer to tokens in the \( u \)-th sequence.
%
Following Proposition \ref{prop-gradient-columns-rows} and \ref{prop-gradient-asymmetric-growth-rows-columns}, we obtain 
%
\begin{equation}
\Delta \bm{w}_{\cdot, k}  =  \sum_{u=1}^D \left[[\bm{y}]_k \sum_{i=k_u+1}^N \beta^u_{i*}\bm{x}^u_i + \sum_{j=1}^{k_u-1} \beta^u_{*j}[\bm{x}^u_j]_k\bm{y}
\right] \,,
\end{equation}
%
while the squared norm of the weight update of the $m$-th row by
%
\begin{equation}
\Delta \bm{w}_{m, \cdot}  =  \sum_{u=1}^D \left[ \sum_{i=k_u+1}^N \beta^u_{i*}[\bm{x}^u_i]_m\bm{y} + [\bm{y}]_m\sum_{j=1}^{k_u-1} \beta^u_{*j}\bm{x}^u_j
\right]\,,
\end{equation}
%
where in both equations the superscript $u$ indicates tokens belonging to the $u$-th sequence in the dataset, and $k_u$ indicates the position index of $t^*$.
%
It follows again from Proposition \ref{prop-gradient-asymmetric-growth-rows-columns} that the expected value of the squared norm of these updates is
%
\begin{equation}
\frac{\mathbb{E}_{\mathcal{D}}\left[ ||\Delta \bm{w}_{\cdot, k}||^2 \right]}{\mathbb{E}_{\mathcal{D}}\left[ ||\Delta \bm{w}_{m, \cdot}||^2 \right]} \approx \frac{\Gamma_{kk}\text{Tr}(\Sigma)\big(\sum_{u=1}^D\sum_{i=k_u+1}^N \beta^u_{i*}\big) + \Sigma_{kk}\text{Tr}(\Gamma)\big(\sum_{u=1}^D\sum_{j=1}^{k_u-1} \beta^u_{*j}\big)}{\Sigma_{mm}\text{Tr}(\Gamma)\big(\sum_{u=1}^D\sum_{i=k_u+1}^N \beta^u_{i*}\big) + \Gamma_{mm}\text{Tr}(\Sigma)\big(\sum_{u=1}^D\sum_{j=1}^{k_u-1} \beta^u_{*j}\big)} \,.
\end{equation}
%
Following Proposition \ref{prop-counting-prediction-context}, let $\mu_c(t^*, U)$ be the random variable quantifying the number of tokens predicted by $t^*$, while $\mu_p(t^*, U)$ is the random variable quantifying the number of tokens predicted by $t^*$.
%
When assuming bidirectional training training, it follows that the empirical average of $\mu_c(t^*, U)$ and $\mu_p(t^*, U)$ over the dataset $\{U_1, \dots, U_D\}$ approaches,
%
\begin{equation}
    \frac{\mathbb{E}_P[\mu_c(t^*, U)]}{\mathbb{E}_P[\mu_p(t^*, U)]} = 1 \,.
\end{equation}
%
From Proposition~\ref{prop-gradient-asymmetric-growth-rows-columns}, it follows that the net increase in the norms of the columns and rows of \( \bm{W}_{qk} \) is the same.
%
Assuming the same initialization conditions as in Theorem~\ref{theo-gradient-directionality}, it follows that by the end of training,
%
\begin{equation}
    \mathrm{Var}(\|\bm{w}_{m, \cdot}\|) = \mathrm{Var}(\|\bm{w}_{\cdot, k}\|) \quad \forall\, k, m \,,
\end{equation}
%
which concludes the proof.
\end{proof}
%
%

We now prove Theorem~\ref{theo-gradients-symmetry}, showing that the bidirectional training objective induces approximate symmetry in the gradient contributions of each token pair \( (i, j) \).
%
Specifically, we assume that predicting token \( t_i \) from \( t_j \) is correlated with predicting \( t_j \) from \( t_i \), that is, the terms \( \beta_{ij} \) and \( \beta_{ji} \) from Proposition~\ref{prop-gradients-self-attention} are correlated, leading to similar contributions in both directions.
%
\begin{proof}
%
%
Let $U = \{t_1, \dots, t_N\}$ be a sequence of tokens.
%
It follows from Proposition \ref{prop-gradients-self-attention} that the implciit weight update for $\bm{W}_{qk}$ following the gradient of $\mathcal{L}(\bm{U})$ w.r.t. $\bm{W}_{qk}$ is given by
%
\begin{equation}
\label{eq:gradient-self-attention-summations}
    \Delta \bm{W}_{qk} = \sum_{i =1}^N \sum_{j=1}^N \beta_{ij} \bm{K}_{ij} \,,
\end{equation}
%
where we neglect any constant of proportionality (e.g. learning rate) for simplicity.
%
The double summation in Equation \eqref{eq:gradient-self-attention-summations} contains $N^2$ elements.
%
Note that \( \bm{K}_{ji} = \bm{K}_{ij}^\top \), so we can rewrite the double summation as follows,
%
\begin{equation}
\sum_{i = 1}^N \sum_{j =1}^N \beta_{ij} \bm{K}_{ij} =
\sum_{i =1}^N \beta_{ii} \bm{K}_{ii} +
\sum_{\substack{i,j = 1 \\ i < j}}^N(\beta_{ij} \bm{K}_{ij} + \beta_{ji} \bm{K}_{ji})
\end{equation}
%
where the first term includes the diagonal terms, and the second includes the contributions of every pair $(i,j)$ with $i,j \in [0, \dots, N]$.
%
The second term can be written as,
%
\begin{equation}
 \sum_{\substack{i,j = 1 \\ i < j}}^N (\beta_{ij} \bm{K}_{ij} + \beta_{ji} \bm{K}_{ji}) = \sum_{\substack{i,j = 1 \\ i < j}}^N(\beta_{ij} \bm{K}_{ij} + \beta_{ji} {\bm{K}_{ij}}^\top) \,, 
\end{equation}
%
and by decomposing $\bm{K}_{ij}$ in its symmetric and skew-symmetric parts, such that $\bm{K}_{ij} = \bm{S}_{ij} + \bm{N}_{ij}$, we obtain,
%
\begin{equation}
\label{eq:gradients-symmetric-terms}
\begin{split}
    \sum_{\substack{i,j = 1 \\ i < j}}^N (\beta_{ij} \bm{K}_{ij} + \beta_{ji} {\bm{K}_{ij}}^\top)
    & = \sum_{\substack{i,j = 1 \\ i < j}}^N \big[ \beta_{ij}(\bm{S}_{ij} + \bm{N}_{ij}) + \beta_{ji}(\bm{S}_{ij} + \bm{N}_{ij})^T \big] \\
    & = \sum_{\substack{i,j = 1 \\ i < j}}^N \big[(\beta_{ij} + \beta_{ji})\bm{S}_{ij} + (\beta_{ij} - \beta_{ji})\bm{N}_{ij}\big]\,.
\end{split}
\end{equation}
%
Let $\beta_{ij}$ and $\beta_{ji}$ be such that $\text{sign}(\beta_{ij}) = \text{sign}(\beta_{ji})$ and $|\beta_{ij}| \approx $$|\beta_{ji}|$.
%
By defining $\Delta \bm{W}_{qk}\big|_{\bm{t}_i \leftrightarrow \bm{t}_j} = \beta_{ij}\bm{K}
_{ij} + \beta_{ji}{\bm{K}_{ij}}^T$,
%
it follows that
%
\begin{equation}
 \Delta \bm{W}_{qk}\big|_{\bm{t}_i \leftrightarrow \bm{t}_j} \approx \sum_{\substack{i,j = 1 \\ i < j}}^N \beta_{ij}\bm{S}_{ij} = \sum_{\substack{i,j = 1 \\ i < j}}^N \beta_{ij}{\bm{S}}^\top_{ij} = \Delta {\bm{W}_{qk}}^\top\big|_{\bm{t}_i \leftrightarrow \bm{t}_j} \,
\end{equation}
%
thus concluding the proof.
%
\end{proof}
%
Encoder-only models are typically not trained to predict every token in a sequence, but rather a random subset of tokens, and the model can attend to tokens bidirectionally.
%
This is usually called Masked Language Modeling (MLM) \citep{devlinBERTPretrainingDeep2019, liuRoBERTaRobustlyOptimized2019, lanALBERTLiteBERT2020, warnerSmarterBetterFaster2024}.
%
Therefore, only a subset of terms in the double summation of Equation \eqref{eq:gradients-symmetric-terms} has the symmetric properties described above. 
%
We generalize the proof to this case in the following.
%
\input{ICML2025/math/statements/remark-gradient-symmetry-MLM}
%
Let $|M| = pN$ with $0<p<1$ being the percentage of tokens to be predicted during bidirectional training.
%
The total number of pairs in the second term of Equation \eqref{eq:factorization-summation} is given by a binomial coefficient, thus the total number of elements in the summation is $pN(pN - 1)$.
%
The total number of elements in the third term is instead the product $pN(N - pN)$.
%
Therefore, the percentage of symmetric weight updates from the second term over the total number of updates in the third term is given by
%
\begin{equation}
    \frac{pN(pN - 1)}{pN(N - pN)} \approx \, \frac{pN}{(N - pN)}\,,
\end{equation}
%
in the limit of large $N$.
%
In practice, $p$ is set to be around 15\%-30\% \citep{devlinBERTPretrainingDeep2019, liuRoBERTaRobustlyOptimized2019,lanALBERTLiteBERT2020,warnerSmarterBetterFaster2024}, leading to $ \approx 25\%$ of symmetric weight updates on average.
%





%% PROOF PROPERTIES OF SYMMETRY SCORE
\subsection{Properties of the symmetry score in Definition \ref{def-symmetry-score} and related proofs}
\label{supp-math-symmetry-score}
%
The score $s$ we introduce in Section~\ref{sec-results-pretrained-models} indicates the \emph{degree} of symmetry of a matrix $\bm{M}$ by quantifying the contribution to the Frobenious norm of its symmetric and skew-symmetric parts. 
%
In particular, $s$ equals 1 and -1 for a fully symmetric and skew-symmetric matrix.
%
Accordingly, positive (negative) values of $s$ indicate the presence of symmetric (skew-symmetric) structures.
%
Here, we provide a proof for these properties. 
%
First, we show that the Frobenious norm of any square matrix $\bm{M}$ can be decomposed in the sum of the Frobenious norm of its symmetric and skew-symmetric components, as in the following Lemma,
%
\input{ICML2025/math/statements/lemma-decomposition-norm-symmetric-skew-symmetric}
\input{ICML2025/math/proofs/proof-lemma-decomposition-norm-symmetric-skew-symmetric}
%
Next, we formulate the properties of the symmetry score as follows,
%
\input{ICML2025/math/statements/prop-score-symmetry-properties}
\input{ICML2025/math/proofs/proof-prop-symmetry-score-properties}
%



%% PROOF PROPERTIES OF DIRECTIONALITY SCORE
\subsection{Properties of the directionality score in Definition \ref{def-directionality-score} and related proofs}
\label{supp-math-directionality-score}
%
The score $d$ we introduce in Section~\ref{sec-results-pretrained-models} quantifies the directional bias of a square matrix $\bm{M}$ by comparing the total norm of the "outliers" rows and columns, that is, that are higher than $\gamma$ times the standard deviations of the norms.
%
A directionality score $d$ of 1 indicates the presence of rows with high ``outlier'' norms and the absence of outliers in the distribution of the column norms.  
%
The opposite is true for a directionality score $d$ of -1.
%
Accordingly, positive (negative) values of $d$ indicate the presence of row (column) dominance in the matrix.
%
Here, we provide a proof for these properties.
%
\input{ICML2025/math/statements/prop-score-directionality-properties}
\input{ICML2025/math/proofs/proof-prop-directionality-score-properties}